\scanpage{5}\mspage{3}
% ll. 1–3
\inserted{Par descente plate,} on peut supposer $X$ local \deleted{complet} \inserted{hensélien}.
\inserted{\unsure{car|on} n'aura pas \unsure{utilisé|seul.}.} \deleted{(gr. \unsure{au} \unsure{m}) \unsure{l'existence} des solutions $G$}
\deleted{Que l'hypothèse sur les Ass est conservée par passage}

% Upper left, composite marginal sketch. The short undirected connectors have no head.
\diagram{\begin{tikzcd}[ampersand replacement=\&,row sep=tiny,column sep=small]
G_{U_n}\& {} \& {} \& {} \& {}\\
U_n\arrow[d,phantom,"\cap" description]\&x_1\arrow[d]\&x'\arrow[l]\arrow[dl]\&\widehat{x}_1\arrow[dr]\arrow[d]\&\widehat{x}'\in\operatorname{Ass}\widehat U\;\deleted{\unsure{\text{générisant}}}\arrow[d]\\
{}\&x\&{}\&\deleted{\unsure{x}}\&\widehat x\\
X_n\arrow[r,dash]\arrow[d,dash,"f_n"]\&X\arrow[d,dash]\&\widehat X=X'\supset U'\arrow[l]\&{}\&{}\\
Y_n\&Y\arrow[l]\&{}\&{}\&{}\\
{}\&y\&y'\arrow[l,dash]\&{}\&{}
\end{tikzcd}}

% ll. 3–8
détruit les hypothèses \inserted{\textcircled{a} \textcircled{b}} faites au départ, \deleted{\unsure{mais}} \inserted{\unsure{ni|en}}
\deleted{(quant $X^h\colon\operatorname{Spec}B^h$) on aura gagné (puisque $B^h$}
l'hypothèse que $B=\mathcal O_{X,x}$ est à fibres formelles
géom\textsuperscript{t}. normales (ce qui est en effet une
condition stable par hensélisation). Par contre,
on aura gagné que le morphisme $X'=\operatorname{Spec}\widehat B\to X=\operatorname{Spec}B$ est
à \uline{fibres formelles \inserted{(géom)} intègres}, donc si $Z_\alpha$ sont les
cycles associés à $G$ sur $X$, \deleted{et \unsure{à} $X'$} alors \deleted{\unsure{el}} ceux
associés à $\deleted{\unsure{F}}G'$ sur $X'$ sont les $Z'_\alpha$ (qui sont,
encore intègres !). Donc l'hyp. \textcircled{a}, \inserted{savoir $Z'_\alpha\cap U_y\ne\varnothing$ pr $\forall\,\alpha$,} reste alors vérifiée
en remplaçant $(X,G,U)$ par $(X',G',U')$. Cela nous
permet de remplacer $X$ par $X'$, (l'hyp. \textcircled{b} reste
trivialement vérifiée).

% ll. 15–end
Utilisant le \unsure{syst.} des solutions
du Pb obtenu pour les $Y_n$, \deleted{et} passant à la
limite, on trouve donc un module quotient
$G\to H$ de $G$, qui est plat sur $Y$, \deleted{\unsure{et}} et qui est
tel que l'\unsure{hom|homo|homom.} $G\to H$ soit \deleted{\unsure{un}} un isom sur les
pts de $U_y$ ; comme pour $\forall\,Z_\alpha$, on a $Z_\alpha\cap U_y\ne\varnothing$,
on conclut que $G\to H$ est injectif sur $U$, \deleted{donc}
injectif puisque $G\to i_*(G_U)$ est injectif. Donc
$G=H$, i.e. $G$ est \uline{plat} sur $Y$, en \unsure{$=$} temps on
voit que $G_y=H_y$ est tel que $x\notin\operatorname{Ass}G_y$, donc
on a gagné.

% Lower left margin, braced together.
\diagram{\begin{tikzcd}[ampersand replacement=\&,row sep=small,column sep=small]
F_n\arrow[r,"u_n"]\arrow[dr,"\text{surj}"']\&i_{n*}(G_{nU})\\
{}\&G_n\arrow[u]\arrow[r,"\unsure{v_n}"']\&H_n
\end{tikzcd}}
\margin{$H_n=\operatorname{Im}u_n=\operatorname{Im}v_n$ est un quotient de $G_n$}
% archive sheet number: 4
